%--------------------------------%
% Author     : Gajanan Choudhary %
% License    : MIT               %
% About file : LaTeX Resume      %
%--------------------------------%

\documentclass[letterpaper,10pt]{article}

%--------------------------------%
% Author     : Gajanan Choudhary %
% License    : MIT               %
% About file : LaTeX Resume      %
%--------------------------------%

\RequirePackage{ifpdf}[2.3]
\ifpdf
  \pdfminorversion=7
\fi

\usepackage{cmap}  % For fixing copy pasting ligatures - ff, fi, fl, ffl, etc.
%Partially deal with overfull hbox / vbox
\usepackage[T1]{fontenc}
\usepackage[final]{microtype}
\usepackage{lmodern}
%\usepackage{ebgaramond}

% Some packages to write mathematics.
%\usepackage{amsmath,amsthm,amsfonts,amscd} 
%\usepackage{eucal} 	 	% Euler fonts
%\usepackage{epsfig}    % Allows inclusion of eps files.
%\usepackage{citesort}   % 

\usepackage{url} 		% Allows good typesetting of web URLs.

\usepackage{latexsym}
%\usepackage{graphics}
%\usepackage{epstopdf}
%\usepackage{algorithmic}
\usepackage{amssymb}
\usepackage{booktabs, makecell}
%\usepackage{bm}
%\usepackage[section]{placeins}
%\usepackage{xfp}
\usepackage{enumitem}  % For indentation of bullets

%%%%%%%%%%%%%%%%%%%%%%%%%%%%%%%%%%%%%%%%%%%%%%%%%%%%%%%%%%%%%%%%%%%%%%
\usepackage[referable]{threeparttablex}
\usepackage{footnote}
%\usepackage{cite} %% Had to comment out to prevent error
%\usepackage{siamcleveref}       %

%----------------------------------%
% Need the following for Wordcloud %
%----------------------------------%
\usepackage{graphicx}         	% Allows inclusion of eps files.
\usepackage{tikz}
\usepackage{tikzscale}
%\usepackage[percent]{overpic}
%\usetikzlibrary{calc}
%\usetikzlibrary{3d}
\usepackage{wrapfig}
\usepackage[utf8]{inputenc}
\usepackage{anyfontsize}
%\usepackage[export]{adjustbox} % Bounding box on figure
\DeclareGraphicsExtensions{.tikz,.pdf} %.eps,.png,.jpg,

\usepackage{contour}
\usepackage[normalem]{ulem}

\newcommand{\recentyear}{1900} 
\usepackage[american]{babel}
\usepackage{csquotes}             % idk, babel/biber like it
\usepackage[style=apa,
            apamaxprtauth=99
            backend=biber,
            sortcites=true,
            sorting=ydmdnt,
            language=american]{biblatex}      % for reference sections
\addbibresource{CV.bib}

%----------------------------------------------------------------%
% Packages borrowed from (MIT licensed) github.com/sb2nov/resume %
%----------------------------------------------------------------%
\usepackage[empty]{fullpage}
\usepackage{titlesec}
\usepackage{marvosym}
%\usepackage[usenames,dvipsnames]{color}
\usepackage{verbatim}
\usepackage{enumitem}
\usepackage[hidelinks]{hyperref}
\usepackage{fancyhdr}
\setlength{\footskip}{3.60004pt}
%\usepackage[english]{babel}
\usepackage{tabularx}
%-----------------------------------------------------------------%


%--------------------------------%
% Author     : Gajanan Choudhary %
% License    : MIT               %
% About file : LaTeX Resume      %
%--------------------------------%

%----------------------------------------------%
% Edit the following -- Put your details below %
%----------------------------------------------%

\def\Author{Gajanan Choudhary}
\def\AuthorEmail{gajananchoudhary91@gmail.com}
\def\AuthorAddress{Austin, TX, USA}
\def\AuthorPhoneText{\smallplus{}1 (512) 657-3030}
\def\AuthorPhoneLink{tel:15126573030}

% You need to escape special characters in website link.
% See: https://stackoverflow.com/questions/2894710/how-to-write-urls-in-latex
\def\AuthorWebsiteText{https://gajanan-choudhary.github.io/}
\def\AuthorWebsiteLink{https://gajanan-choudhary.github.io/}
\def\AuthorGitHubText{https://github.com/gajanan-choudhary/}
\def\AuthorGitHubLink{https://github.com/gajanan-choudhary/}
\def\AuthorLinkedInText{https://www.linkedin.com/in/gajananchoudhary}
\def\AuthorLinkedInLink{https://www.linkedin.com/in/gajananchoudhary}

%%%----------------------------------------------------------%%%
%%%----------------------------------------------------------%%%
%%% Most likely, you will not need to edit things below this %%%
%%%----------------------------------------------------------%%%
%%%----------------------------------------------------------%%%

%------------------%
% Utility commands %
%------------------%

%---------------------%
% Use sans-serif font %
%---------------------%
%\renewcommand{\familydefault}{\sfdefault}

%--------------%
% Line spacing %
%--------------%
\linespread{1.2}

%-------%
% Dates %
%-------%
\newcommand{\dates}[2]{#1 -- #2}

%---------------------------%
% For underlining hyperlink %
%---------------------------%
\renewcommand{\ULdepth}{1.8pt}
\contourlength{0.8pt}
\newcommand{\myuline}[1]{%
  \uline{\phantom{#1}}%
  \llap{\contour{white}{#1}}%
}

%-----------%
% Wordcloud %
%-----------%
\definecolor{uto}{RGB}{191, 87, 0}
\definecolor{utg}{RGB}{51, 63, 72}

% The following is needed for properly scaling text in the wordcloud. See:
% https://tex.stackexchange.com/questions/328228/appearance-of-tiny-or-scriptsize-fontsize-in-latex-horizontal-stretch/328237
%%% TEXT FONT FIX
\rmfamily
\DeclareFontShape{T1}{lmr}{bx}{sc}{<-> cmr10}{}% USE BOLD SCSHAPE NOT OTHERWISE DEFINED
%%% MATH FONT FIX
\DeclareFontFamily{OML}{zlmm}{}
\DeclareFontShape{OML}{zlmm}{m}{it}{<-> lmmi10}{}
\DeclareFontShape{OML}{zlmm}{b}{it}{<->ssub * zlmm/m/it}{}
\DeclareFontShape{OML}{zlmm}{bx}{it}{<->ssub * zlmm/m/it}{}

\DeclareMathVersion{Tinyb}
\SetSymbolFont{operators}{Tinyb}{T1}{lmr}{bx}{sc}
\SetSymbolFont{letters}{Tinyb}{OML}{zlmm}{m}{it}
%%%
\newenvironment{tinyb}{\bgroup\tiny\bfseries\scshape\mathversion{Tinyb}}{\egroup} % edit: now with \tiny included


%---------------%
% Writing "C++" %
%---------------%
% See: https://tex.stackexchange.com/questions/4302/prettiest-way-to-typeset-c-cplusplus
\def\CC{{C\nolinebreak[4]\hspace{-.05em}\raisebox{.4ex}{\tiny\bf ++}}}
% Needed for wordcloud:
\def\bigCC{{C\nolinebreak[4]\hspace{-.05em}\raisebox{.4ex}{\normalsize\bf ++}}}

%---------------------%
% Writing smaller "+" %
%---------------------%
\def\smallplus{\nolinebreak[2]\raisebox{.4ex}{\tiny\bf +}}

%---------------------------------------------------------------------------%
% Utility commands -- Borrowed from (MIT licensed) github.com/sb2nov/resume %
% Note: I have modified a few things after borrowing.                       %
%---------------------------------------------------------------------------%

\pagestyle{fancy}
\fancyhf{} % clear all header and footer fields
\fancyfoot{}
\renewcommand{\headrulewidth}{0pt}
\renewcommand{\footrulewidth}{0pt}

% Adjust margins
\addtolength{\oddsidemargin}{-0.25in}
\addtolength{\evensidemargin}{-0.25in}
\addtolength{\textwidth}{0.5in}
\addtolength{\topmargin}{-0.15in}
\addtolength{\textheight}{0.5in}

\urlstyle{same}

\raggedbottom
\raggedright
\setlength{\tabcolsep}{0in}

%---------------------------------------------------------------------------%
% Sections formatting
\titleformat{\section}{
  \vspace{-19pt}\scshape\raggedright\large
}{}{0em}{\vspace{-4pt}\color{gray}\titlerule[0.5pt]\color{black}\\}[\vspace{-6pt}]
%}{}{0em}{\color{gray}\rule{\linewidth}{0.5pt}\color{black}\\}[\vspace{-6pt}]


%---------------------------------------------------------------------------%
%---- gkc Modifying ----%
%\newcommand{\resumeSubheading}[4]{  % Original
\newcommand{\resumeSubheading}[2]{
  \vspace{-4pt}\item
    \begin{tabularx}{\linewidth}[t]{Xr}
      \textbf{#1} & #2 \\
    \end{tabularx}\vspace{-13pt}
}

\newcommand{\resumeSubSubheading}[2]{
    \vspace{-3pt}
    \begin{tabularx}{\linewidth}{Xr}
      #1 & #2 \\
    \end{tabularx}\vspace{-13pt}
}

%---------------------------------------------------------------------------%
\newcommand{\resumeItem}[2]{
  \item{
    {\textit{#1}}{: #2 \vspace{-2.5pt}}
  }
}

\newcommand{\resumeSubItem}[2]{
  \item{

    \textbf{#1}{: #2 \vspace{-7pt}} 
%    \textbf{#1}{: #2 \vspace{-9pt}}  % Use this if 11pt article
  }
}

\newcommand{\resumeUnlabeledItem}[1]{\item{{#1 \vspace{-2.5pt}}}}

%\newcommand{\resumeSubItem}[2]{\resumeItem{#1}{#2}\vspace{-4pt}}
%---------------------------------------------------------------------------%
% Bullet size and type
\renewcommand{\labelitemi}{\raisebox{1pt}{\footnotesize{$\bullet$}}}
\renewcommand{\labelitemii}{\raisebox{1pt}{\footnotesize{$\circ$}}}

%---------------------------------------------------------------------------%
\newcommand{\resumeSubHeadingListStart}{\begin{itemize}[leftmargin=0.225in]}
\newcommand{\resumeSubHeadingListEnd}{\end{itemize}\vspace{-3pt}}

%---------------------------------------------------------------------------%
\newcommand{\resumeSubItemListStart}{\vspace{-2pt}\begin{itemize}[leftmargin=0.225in]}
\newcommand{\resumeSubItemListEnd}{\end{itemize}\vspace{2pt}}

%---------------------------------------------------------------------------%
\newcommand{\resumeItemListStart}{\vspace{-7pt}\begin{itemize}[leftmargin=0.225in]}
\newcommand{\resumeItemListEnd}{\end{itemize}}

%---------------------------------------------------------------------------%
% Borrowed from:
%     https://github.com/btskinner/tex_cv
% If you want to go the partially manual route, look at another repo:
%     https://github.com/rgeirhos/academic-cv-publications

% ------------------------%
% BibLaTeX sorting scheme %
% ------------------------%

% Create a new sorting scheme that will sort in reverse chronological
% order, taking months into account correctly (i.e., in chronological
% order and not alphabetical order)

\DeclareSortingTemplate{ydmdnt}{
  \sort{
    \field{presort}
  }
  \sort[final]{
    \field{sortkey}
  }
  \sort[direction=descending]{
    \field{sortyear}
    \field{year}
    \literal{9999}
  }
  \sort[direction=descending]{
    \field[padside=left,padwidth=2,padchar=0]{month}
    \literal{99}
  }
  \sort[direction=descending]{
    \field[padside=left,padwidth=2,padchar=0]{day}
    \literal{99}
  }
  \sort{
    \field{sortname}
    \field{author}
    \field{editor}
    \field{translator}
    \field{sorttitle}
    \field{title}
  }
  \sort{
    \field{sorttitle}
  }
  \sort[direction=descending]{
    \field[padside=left,padwidth=4,padchar=0]{volume}
    \literal{9999}
  }
}

%------------------%
% BibLaTeX options %
%------------------%

% language mapping
\DeclareLanguageMapping{american}{american-apa}

% ignore addendum and note fields that are sometimes useful, but not here
\DeclareSourcemap{
  \maps[datatype=bibtex]{
    \map{
      \step[fieldset=addendum, null]
      \step[fieldset=note, null]
    } 
  }
}

% macro so titles are converted to doi, url, or isbn link (if available)
% h\t https://tex.stackexchange.com/a/48409
\ExecuteBibliographyOptions{doi=false,url=false,isbn=false}
\newbibmacro{string+doiurlisbn}[1]{%
  \iffieldundef{doi}{%
    \iffieldundef{url}{%
      \iffieldundef{isbn}{%
        \iffieldundef{issn}{%
          #1%
        }{%
          \href{http://books.google.com/books?vid=ISSN\thefield{issn}}{#1}%
        }%
      }{%
        \href{http://books.google.com/books?vid=ISBN\thefield{isbn}}{#1}%
      }%
    }{%
      \href{\thefield{url}}{#1}%
    }%
  }{%
    \href{http://dx.doi.org/\thefield{doi}}{#1}%
  }%
}

\DeclareFieldFormat{title}{\usebibmacro{string+doiurlisbn}{\mkbibemph{#1}}}
\DeclareFieldFormat[article,incollection]{title}%
{\usebibmacro{string+doiurlisbn}{\mkbibquote{#1}}}

% no letter for same year entries
\defbibenvironment{mybib}
{\list
  {}
  {\setlength{\leftmargin}{\bibhang}%
    \setlength{\itemindent}{-\leftmargin}%
    \setlength{\itemsep}{\bibitemsep}%
    \setlength{\parsep}{\bibparsep}}}
{\endlist}
{\clearfield{extradate}\item}

% filter by year (only print those after year value in \recentyear
% macro above)
\defbibcheck{recent}{
  \iffieldint{year}
  {\ifnumless{\thefield{year}}{\recentyear}
    {\skipentry}
    {}}
  {\skipentry}
}



\begin{document}

%-----------------------------%
% Customize your resume below %
%-----------------------------%
\href{\AuthorWebsiteLink}{\scshape\huge\Author}\vspace{2pt}
%\href{\AuthorWebsiteLink}{\scshape\huge\Author}\vspace{2pt}

\begin{small}
  %Website: \href{\AuthorWebsiteLink}{\color[HTML]{0f4cb4}\myuline{\smash{\AuthorWebsiteText}}}
  %Website: \href{\AuthorWebsiteLink}{\AuthorWebsiteText}
  LinkedIn: \href{\AuthorLinkedInLink}{\AuthorLinkedInText}

  \vspace{2pt}
  GitHub: \href{\AuthorGitHubLink}{\AuthorGitHubText}

  \vspace{1pt}
  \AuthorAddress{} |
  \href{mailto:\AuthorEmail}{\AuthorEmail} |
  \href{\AuthorPhoneLink}{\AuthorPhoneText}
\end{small}

%------------------------------------SUMMARY-----------------------------------%
\section{Summary}
Performance engineer with 13 years of interdisciplinary research,
artificial intelligence (AI), and high-performance computing (HPC) production
software development experience spanning CPU/GPU/custom AI accelerator kernel
libraries, scientific computing, deep learning, and applied mathematics.
Creator of and contributor to numerous AI/HPC software written in
C/C++, SYCL, Python, and Fortran, and author of 6 technical publications.

%------------------------------PROGRAMMING SKILLS------------------------------%
\section{Skills}
  \resumeSubItemListStart
    \resumeSubItem{Programming}{C/\CC{}, SYCL, Python, Fortran, Assembly,
                               MPI, OpenMP, Bash, and Python/C API.}
    \resumeSubItem{Tools}{Git, GitHub, Cursor, Claude Code, Actions, Docker,
                         CMake, Make, Gcov, LCOV, Doxygen, and \LaTeX{}.}
  \resumeSubItemListEnd

%------------------------------------HEADING-----------------------------------%
%\setlength{\columnsep}{0pt}%
\vspace{-241.5pt}
%%\vspace{-162pt}
%\vspace{-82pt}
%\vspace{-168.25pt}

\begin{wrapfigure}[-20]{R}{0pt} %{-3cm}

  \centering

  \begin{tikzpicture}[x=1in, y=-1in, every node/.style={right}] 

    \begin{tinyb}

      \node[color=utg]at (0.12,0.12){\footnotesize x86 Assembly\space };

      \node[color=uto]at (0.16,0.20){\tiny Kernel optimizations\space };

      \node[color=uto]at (0.57,0.26){\tiny Docker\space };

      \node[color=utg]at (0.17,0.34){\Large Python\space };

      \node[color=uto]at (0.10,0.44){\tiny Cursor\space };

      \node[color=uto]at (0.36,0.48){\footnotesize Sparse BLAS\space };

      \node[color=utg]at (0.12,0.56){\scriptsize SWIG\space };

      \node[color=utg]at (0.43,0.61){\normalsize MPI\space };

      \node[color=utg]at (0.73,0.62){\tiny Algebra\space };

      \node[color=uto]at (0.1,0.85){\tiny MoE\space };

      \node[color=utg]at (0.30,0.86){\scriptsize GEMMs\space };

      \node[color=uto]at (0.16,0.96){\scriptsize Intel Corporation\space };

      \node[color=utg]at (0.13,1.06){\scriptsize PyTorch\space };

      \node[color=uto]at (0.16,1.18){\small CMake\space };

      \node[color=uto]at (0.83,0.17){\large SYCL\space };

      \node[color=utg]at (0.81,0.32){\scriptsize GPUs\space };

      \node[color=uto]at (0.96,0.58){\footnotesize ctypes\space };

      \node[color=utg]at (0.85,0.89){\footnotesize IIT Kharagpur\space };

      \node[color=utg  ]at (1.11,0.98){\scriptsize Bash\space };

      \node[color=uto]at (0.51,1.07){\normalsize Machine learning\space };

      \node[color=utg]at (0.65,1.17){\footnotesize Neural Networks\space };

      \node[color=uto]at (1.28,0.14){\scriptsize Performance engineering\space };

      \node[color=uto]at (1.29,0.21){\tiny Claude\space };

      \node[color=utg]at (1.52,0.22){\tiny Image Processing\space };

      \node[color=utg]at (1.66,0.32){\footnotesize CI/CD\space };

      \node[color=utg]at (1.70,0.40){\tiny Benchmarking\space };

      \node[color=uto]at (1.07,0.35){\large GitHub\space };

      \node[color=utg]at (1.06,0.48){\footnotesize HPCG\space };

      \node[color=uto]at (1.41,0.47){\scriptsize Accelerators\space };

      \node[color=uto]at (1.44,0.56){\tiny HW-SW Co-design\space };

      \node[color=utg]at (1.36,0.62){\small Fortran\space };

      \node[color=utg]at (1.60,0.87){\tiny CPUs\space };

      \node[color=uto]at (1.80,0.90){\footnotesize AdH\space };

      \node[color=uto]at (1.46,0.96){\tiny MATLAB\space };

      \node[color=uto]at (1.57,1.05){\scriptsize Tenstorrent\space };

      \node[color=utg]at (1.54,1.19){\scriptsize f2py\space };

      \node[color=uto]at (1.77,1.17){\large CUDA\space };

      \node[color=utg]at (2.39,0.12){\footnotesize OpenMP\space };

      \node[color=uto]at (2.42,0.22){\scriptsize UT Austin\space };

      \node[color=uto]at (2.08,0.23){\scriptsize Make\space };

      \node[color=utg]at (2.16,0.32){\small x86 Intrinsics\space };

      \node[color=utg]at (2.33,0.41){\tiny Scaling\space };

      \node[color=utg]at (2.73,0.42){\tiny GPT\space };

      \node[color=uto]at (1.94,0.50){\normalsize Applied Math\space };

      \node[color=utg]at (2.61,0.58){\tiny Libraries\space };

      \node[color=utg]at (2.26,0.61){\scriptsize ERDC\space };

      \node[color=utg]at (2.12,0.86){\tiny ANSYS\space };

      \node[color=uto]at (2.45,0.87){\tiny JIT\space };

      \node[color=uto]at (2.65,0.90){\footnotesize Git\space };

      \node[color=utg]at (2.05,1.03){\huge C/\bigCC{}\space };

      \node[color=uto]at (2.33,1.13){\tiny Scale-up\space };

      \node[color=utg]at (2.40,1.19){\tiny Scale-out\space };

      \node[color=utg]at (2.70,1.12){\scriptsize AI\space };

      \node[color=utg]at (2.74,1.20){\tiny LLMs\space };

    \end{tinyb}

    \node[color=uto]at (0.10,0.74){\huge Computational Scientist\space };

    \draw[color=gray] (current bounding box.north east) -- (current bounding box.north west) -- (current bounding box.south west) -- (current bounding box.south east) -- cycle;

  \end{tikzpicture}

\end{wrapfigure}

\strut{} %% For ending wrapfig wrapping properly

\vspace{225.65pt}

% Following needed since wordcloud messes up wrapping the \hrule lines

%%\parbox{0.99\linewidth}{
%  %\begin{center}
%      \href{\AuthorWebsiteLink}{\scshape\huge\Author}\vspace{2pt}
%  %\end{center}
%%}% End of: \parbox{0.546\linewidth}{
%
%  %\begin{center}
%    \begin{small}
%      %Website: \href{\AuthorWebsiteLink}{\color[HTML]{0f4cb4}\myuline{\smash{\AuthorWebsiteText}}}
%      Website: \href{\AuthorWebsiteLink}{\AuthorWebsiteText}
%
%      \vspace{2pt}\AuthorAddress
%
%      \vspace{2pt}\href{mailto:\AuthorEmail}{\AuthorEmail} |
%      \href{\AuthorPhoneLink}{\AuthorPhoneText}
%
%    \end{small}
%  %\end{center}


\vspace{-1pt}
%----------------------------------EXPERIENCE----------------------------------%
\section{Select Experience}
  \resumeSubHeadingListStart

    %---------------------------------Title------------------------------------%
    \resumeSubheading{Senior Staff Software Engineer}{\dates{January 2026}{Present}}

      %-------------------------------Job 6------------------------------------%
      \resumeSubSubheading{\textbf{Tenstorrent}}{Jersey City, NJ}
        \resumeItemListStart
          \resumeItem{Responsibilities}
            {Optimized AI kernels and modules in TT-Metalium\texttrademark,
            Tenstorrent's open-source, low-level AI hardware SDK, to enable
            running frontier LLMs on Tenstorrent AI hardware.}
          \resumeItem{Key Contributions}
            {Delivered a generalized mixture of experts (MoE) module to support
            a swath of LLMs and refactored the codebase to enable
            next-generation, non-backwards-compatible AI accelerator.}
        \resumeItemListEnd

    %---------------------------------Title------------------------------------%
    \resumeSubheading{Math Algorithm Engineer}{\dates{February 2021}{January 2026}}

      %-------------------------------Job 5------------------------------------%
      \resumeSubSubheading{\textbf{Intel Corporation}}{Austin, TX}
        \resumeItemListStart
          \resumeItem{Responsibilities}
            {Owned sparse linear algebra components as a
            contributor in the
            Intel\textsuperscript{\scriptsize{\textregistered}} oneAPI Math
            Kernel Library (oneMKL) team.}
          \resumeItem{Key contributions}
            {Led cross-organizational collaborations for Intel's open-source
            efforts and delivered extensive foundational, architectural,
            performance, and quality improvements on the sparse BLAS domain for
            HPC and AI through math APIs such as matrix-vector and matrix-matrix
            multiplication.}
        \resumeItemListEnd

    %---------------------------------Title------------------------------------%
    % \resumeSubheading{Research Associate}{\dates{December 2020}{February 2021}}
 
    %   %-------------------------------Job 4------------------------------------%
    %   \resumeSubSubheading{\textbf{The University of Texas at Austin}}{Austin, TX}
    %     \resumeItemListStart
    %       \resumeItem{Responsibilities}
    %         {Conducted research on hurricanes and managed research direction
    %         of post-doctoral fellows and doctoral students in the team.}
    %       \resumeItem{Key contribution}
    %         {Led the development of coupled PyTorch neural network and CFD
    %         models in the \textbf{Python} software, Water Coupler, to
    %         potentially save lives and billions of dollars through hurricane
    %         flood forecasts.}
    %     \resumeItemListEnd

    %---------------------------------Title------------------------------------%
    \resumeSubheading{Postdoctoral Fellow}{\dates{October 2019}{December 2020}}
 
      %-------------------------------Job 3------------------------------------%
      \resumeSubSubheading{\textbf{The University of Texas at Austin}}{Austin, TX}
        \resumeItemListStart
          \resumeItem{Responsibilities}
            {Open-source parallel software development for research on coupled
            fluid dynamics models.}

          \resumeItem{Key contributions}
            {Coupled diffusive wave and groundwater models in the \textbf{C}
            software, Adaptive Hydraulics (AdH), created pyADCIRC, a
            \textbf{Python} interface for the \textbf{Fortran} software,
            ADvanced CIRCulation (ADCIRC), and added coupling support for
            pyADCIRC in Water Coupler.}
        \resumeItemListEnd

    %---------------------------------Title------------------------------------%
    \resumeSubheading{Graduate Research Assistant}{\dates{September 2014}{October 2019}}
 
    %---------------------------------Job 2------------------------------------%
      \resumeSubSubheading{\textbf{The University of Texas at Austin}}{Austin, TX}
        \resumeItemListStart
          \resumeItem{Responsibilities}
            {Software and library development for research on coupled
            fluid dynamics models.}

          \resumeItem{Key contributions}
            {Created \textbf{Water Coupler}, an \textbf{open-source Python}
            software for coupling multiple CFD and machine learning models
            written in \textbf{C}, \textbf{\CC{}}, \textbf{Fortran} and
            \textbf{Python} through portable parallel Python interfaces for
            simulating compound floods due to hurricanes.}
       \resumeItemListEnd

%    %---------------------------------Title------------------------------------%
%    \resumeSubheading{Assistant Surveyor}{\dates{July 2013}{July 2014}}
% 
%      %-------------------------------Job 1------------------------------------%
%      \resumeSubSubheading{\textbf{Indian Register of Shipping}}{Mumbai, India}
%        \resumeItemListStart
%          \resumeItem{Responsibilities}
%            {Research on stress response of ship hulls to bending, shear,
%            torsion, and warping loads.}
%
%          \resumeItem{Key contribution}
%            {Developed a \textbf{MATLAB} software with a \textbf{GUI} for
%            modeling and stress analysis of ships.}
%        \resumeItemListEnd
  \resumeSubHeadingListEnd


%-----------------------------------EDUCATION----------------------------------%
\section{Education}
  \resumeSubHeadingListStart

    \resumeSubheading{The University of Texas at Austin}{Austin, TX}

      \resumeSubSubheading{Doctorate (Ph.D.) in Engineering Mechanics}{\dates{August 2014}{December 2019}}

  \resumeSubHeadingListEnd

%------------------------------------THE END-----------------------------------%
\end{document}
